% Wspólna część topologiczna: kohomologia, forma przecięć,
% liczby Bettiego, charakterystyka Eulera i sygnatura.
% Dwie historyczne etykiety identyfikują jedno twierdzenie.
\ifcsname biblioteka@rozdmuchanie@kohomologia@loaded\endcsname
Kohomologię rozdmuchania punktu obliczono wcześniej
w twierdzeniu~\ref{thm:xvii-blowup}.
\else
\expandafter\gdef\csname biblioteka@rozdmuchanie@kohomologia@loaded\endcsname{}
\begin{twierdzenie}[Kohomologia i forma przecięć rozdmuchania punktu]
\label{tw:ga10-blowup}\label{thm:xvii-blowup}
Niech $S$ będzie gładką, zwartą, spójną zespoloną powierzchnią,
$p\in S$, $\pi:\widetilde S=\operatorname{Bl}_pS\to S$ rozdmuchaniem,
$E=\pi^{-1}(p)$ oraz $e=\operatorname{PD}_{\widetilde S}[E]\in
H^2(\widetilde S;\mathbb Z)$ (symbol $\operatorname{PD}$ oznacza dualność
Poincarégo). Wówczas
\[
H^2(\widetilde S;\mathbb Z)=\pi^*H^2(S;\mathbb Z)\oplus\mathbb Z e.
\]
Dla każdej $\alpha\in H^2(S;\mathbb Z)$ zachodzi
$\langle \pi^*\alpha\smile e,[\widetilde S]\rangle=0$ i
$\langle e^2,[\widetilde S]\rangle=-1$.
W innych stopniach $\pi^*:H^k(S;\mathbb Z)\to
H^k(\widetilde S;\mathbb Z)$ jest izomorfizmem.
Przez $\sigma(M)$ rozumiemy sygnaturę rzeczywistej formy
przecięć na $H^2(M;\mathbb R)$; $b_2$ oznacza drugą liczbę
Bettiego, a $\chi_{\mathrm{top}}$ charakterystykę Eulera.
W szczególności
\[
 b_2(\widetilde S)=b_2(S)+1,\qquad
 \chi_{\mathrm{top}}(\widetilde S)=\chi_{\mathrm{top}}(S)+1,
 \qquad \sigma(\widetilde S)=\sigma(S)-1.
\]
\end{twierdzenie}
\begin{proof}
Wybierzmy małą zamkniętą kulę $B\subset S$ o środku $p$,
$U=S\setminus\operatorname{int}B$, i niech
$\widetilde B=\pi^{-1}(B)$. Przestrzeń $\widetilde B$ jest
otoczeniem dyskowym $E\simeq\mathbb P^1(\mathbb C)\simeq S^2$
i deformacyjnie retraktuje się na $E$. Brzegi $B$ oraz
$\widetilde B$ są identyfikowane przez $\pi$ ze sferą $S^3$.
Zatem $S=U\cup B$ i $\widetilde S=U\cup\widetilde B$
mają na przecięciach przestrzeń homotopijnie równoważną $S^3$.

Z ciągów Mayera--Vietorisa (ze współczynnikami $\mathbb Z$)
i równości $H^1(S^3)=H^2(S^3)=0$ dostajemy
\[
 H^2(S;\mathbb Z)\simeq H^2(U;\mathbb Z),\qquad
 H^2(\widetilde S;\mathbb Z)
 \simeq H^2(U;\mathbb Z)\oplus H^2(E;\mathbb Z).
\]
Pierwsza składowa ostatniego izomorfizmu jest obrazem $\pi^*$,
ponieważ $\pi|_{\widetilde B}$ ma obraz w ściągalnej kuli $B$.
Ograniczenie $e$ do $E$ jest klasą
$c_1(N_{E/\widetilde S})=-\eta$, gdzie
$\eta$ jest dodatnim generatorem $H^2(E;\mathbb Z)$;
wynika to z lematu~\ref{lem:xvii-exceptional}.
Dlatego $e$ generuje dodatkowy składnik całkowity.

Dla jasności sprawdzamy także stopień trzeci.
Z dualności Poincarégo--Lefschetza lub z długiego ciągu pary
$(U,\partial U)$ wynika, że homomorfizm łączący
$H^3(\partial U;\mathbb Z)\to H^4(U,\partial U;\mathbb Z)$
jest izomorfizmem: generator sfery granicznej przechodzi
na względną klasę orientacji spójnej rozmaitości $U$.
Stąd $H^3(U;\mathbb Z)\to H^3(\partial U;\mathbb Z)$
jest zerowy. Porównanie obu ciągów Mayera--Vietorisa
pokazuje, że $\pi^*$ jest izomorfizmem również w stopniu trzecim.
Stopnie zero i jeden porównujemy bezpośrednio z ciągów,
a w stopniu czwartym $\pi$ ma stopień $1$, więc
$\pi^*:H^4(S;\mathbb Z)\to H^4(\widetilde S;\mathbb Z)$
jest izomorfizmem (obie grupy są izomorficzne z $\mathbb Z$).

Ortogonalność wynika ze wzoru projekcji i tego, że
$\pi|_E$ jest odwzorowaniem stałym:
\[
 \langle\pi^*\alpha\smile e,[\widetilde S]\rangle
 =\langle (\pi|_E)^*\alpha,[E]\rangle=0.
\]
Samoprzecięcie $e^2=-1$ daje wspomniany lemat.
Ponieważ $\pi$ ma stopień $1$, forma przecięć na
$\pi^*H^2(S;\mathbb R)$ jest identyczna z formą na $S$.
Otrzymujemy więc rozkład ortogonalny
$Q_{\widetilde S}\cong Q_S\oplus(-1)$ na wolnych częściach
kohomologii. To daje $\sigma(\widetilde S)=\sigma(S)-1$.
Zmiana wyłącznie drugiej liczby Bettiego powoduje
$\chi_{\mathrm{top}}(\widetilde S)=\chi_{\mathrm{top}}(S)+1$.
\end{proof}
\fi
